% META: {"id":"TCOMPL-LOG-EX-001","chapitre":"TCOMPL-LOGARITHME-HISTORIQUE","type_objet":"exercice","capacites":["TCOMPL-LOG-C4"],"capacites_codes":["C4"],"methodes":[],"parcours":1,"competences":["calculer"],"duree_min":10,"mode_creation":"ex_nihilo","sources_inspiration":[],"fichier_tex":"chapitres/TCOMPL-LOGARITHME-HISTORIQUE/exercices/TCOMPL-LOG-EX-001.tex","corrige_tex":"chapitres/TCOMPL-LOGARITHME-HISTORIQUE/corriges/TCOMPL-LOG-CO-001.tex","status":"approved","origin":"generated","status_history":["generated","audited","accepted","approved"],"release_acceptance":"RELEASE_OWNER_BATCH_ACCEPTANCE","acceptance_closure_digest":"sha256:9b3ccf9a81c5520fbb7e03b7b2d3e7bf2057a02834b6d908bf3be5f49f3b553f","acceptance_receipt_digest":"sha256:4fad86f0282e0fa4e0419cca1ad6379ae7af5a280c04a4a90880f90a1c044242"}
% BEGIN-VERIFY
% from sympy import ln, simplify, Rational
% A = ln(12) - ln(3) - ln(2)
% assert simplify(A - ln(2)) == 0
% B = ln(Rational(1,5)) + ln(5)
% assert simplify(B) == 0
% END-VERIFY
\begin{exercice}{TCOMPL-LOG-EX-001}{1}{10}
En utilisant l'équation fonctionnelle, exprimer sans calculatrice, sous la forme $k\ln(2)$ ou $0$ :
\begin{enumerate}
  \item $A = \ln(12) - \ln(3) - \ln(2)$ ;
  \item $B = \ln\!\left(\dfrac{1}{5}\right) + \ln(5)$.
\end{enumerate}
\end{exercice}
